\documentclass[UTF8,12pt]{ctexart}
\usepackage[a4paper,margin=24mm]{geometry}
\usepackage{amsmath,amssymb,bm,graphicx,adjustbox}
\newcommand{\fitmath}[1]{\begin{adjustbox}{max width=\linewidth}$\displaystyle #1$\end{adjustbox}}
\setlength{\parindent}{0pt}
\setlength{\parskip}{.7em}
\sloppy
\allowdisplaybreaks
\begin{document}
\section*{2003 年考研数学三 · 真题}
\subsection*{2003年全国硕士研究生招生考试试题}

\subsection*{一、填空题（本题共6小题，每小题4分，满分24分）}

（1）设 \(\displaystyle f(x)=\begin{cases}x^\lambda\cos\dfrac{\displaystyle 1}{\displaystyle x},\allowbreak &x\ne0,\allowbreak \\0,\allowbreak &x=0,\allowbreak \end{cases}\) 其导函数在 \(\displaystyle x=0\) 处连续，则 \(\displaystyle \lambda\) 的取值范围是\_\_\_\_\_\_。


（2）已知曲线 \(\displaystyle y=x^3-3a^2x+b\) 与 \(\displaystyle x\) 轴相切，则 \(\displaystyle b^2\) 可以通过 \(\displaystyle a\) 表示为 \(\displaystyle b^2=\)\_\_\_\_\_\_。


（3）设 \(\displaystyle a>0\)，\(\displaystyle f(x)=g(x)=\begin{cases}a,\allowbreak &0\le x\le1,\allowbreak \\0,\allowbreak &\text{其他},\allowbreak \end{cases}\) 而 \(\displaystyle D\) 表示全平面，则 \(\displaystyle I=\displaystyle \iint_Df(x)g(y-x)\,\mathrm dx\mathrm dy=\)\_\_\_\_\_\_。


（4）设 \(\displaystyle n\) 维向量 \(\displaystyle \alpha=(a,\allowbreak 0,\allowbreak \ldots,\allowbreak 0,\allowbreak a)^{\mathrm T}\)，\(\displaystyle a<0\)；\(\displaystyle E\) 为 \(\displaystyle n\) 阶单位矩阵，矩阵 \(\displaystyle A=E-\alpha\alpha^{\mathrm T}\)，\(\displaystyle B=E+\dfrac{\displaystyle 1}{\displaystyle a}\alpha\alpha^{\mathrm T}\)，其中 \(\displaystyle A\) 的逆矩阵为 \(\displaystyle B\)，则 \(\displaystyle a=\)\_\_\_\_\_\_。


（5）设随机变量 \(\displaystyle X\) 和 \(\displaystyle Y\) 的相关系数为 \(\displaystyle 0.9\)，若 \(\displaystyle Z=X-0.4\)，则 \(\displaystyle Y\) 与 \(\displaystyle Z\) 的相关系数为\_\_\_\_\_\_。


（6）设总体 \(\displaystyle X\) 服从参数为2的指数分布，\(\displaystyle X_1,\allowbreak X_2,\allowbreak \ldots,\allowbreak X_n\) 为来自总体 \(\displaystyle X\) 的简单随机样本，则当 \(\displaystyle n\to\infty\) 时，\(\displaystyle Y_n=\dfrac{\displaystyle 1}{\displaystyle n}\displaystyle \sum_{i=1}^nX_i^2\) 依概率收敛于\_\_\_\_\_\_。


\subsection*{二、选择题（本题共6小题，每小题4分，满分24分）}

（1）设 \(\displaystyle f(x)\) 为不恒等于零的奇函数，且 \(\displaystyle f^{\prime}(0)\) 存在，则函数 \(\displaystyle g(x)=\dfrac{\displaystyle f(x)}{\displaystyle x}\)（　）。
（A）在 \(\displaystyle x=0\) 处左极限不存在。（B）有跳跃间断点 \(\displaystyle x=0\)。（C）在 \(\displaystyle x=0\) 处右极限不存在。（D）有可去间断点 \(\displaystyle x=0\)。


（2）设可微函数 \(\displaystyle f(x,\allowbreak y)\) 在点 \(\displaystyle (x_0,\allowbreak y_0)\) 取得极小值，则下列结论正确的是（　）。
（A）\(\displaystyle f(x_0,\allowbreak y)\) 在 \(\displaystyle y=y_0\) 处的导数等于零。（B）\(\displaystyle f(x_0,\allowbreak y)\) 在 \(\displaystyle y=y_0\) 处的导数大于零。（C）\(\displaystyle f(x_0,\allowbreak y)\) 在 \(\displaystyle y=y_0\) 处的导数小于零。（D）\(\displaystyle f(x_0,\allowbreak y)\) 在 \(\displaystyle y=y_0\) 处的导数不存在。


（3）设 \(\displaystyle p_n=\dfrac{\displaystyle a_n+|a_n|}{\displaystyle 2}\)，\(\displaystyle q_n=\dfrac{\displaystyle a_n-|a_n|}{\displaystyle 2}\)，\(\displaystyle n=1,\allowbreak 2,\allowbreak \cdots\)，则下列命题正确的是（　）。
（A）若 \(\displaystyle \sum_{n=1}^{\infty}a_n\) 条件收敛，则 \(\displaystyle \sum p_n\) 与 \(\displaystyle \sum q_n\) 都收敛。（B）若 \(\displaystyle \sum_{n=1}^{\infty}a_n\) 绝对收敛，则 \(\displaystyle \sum p_n\) 与 \(\displaystyle \sum q_n\) 都收敛。（C）若 \(\displaystyle \sum a_n\) 条件收敛，则 \(\displaystyle \sum p_n\) 与 \(\displaystyle \sum q_n\) 的敛散性都不定。（D）若 \(\displaystyle \sum a_n\) 绝对收敛，则 \(\displaystyle \sum p_n\) 与 \(\displaystyle \sum q_n\) 的敛散性都不定。


（4）设3阶矩阵 \(\displaystyle A=\begin{pmatrix}a&b&b\\b&a&b\\b&b&a\end{pmatrix}\)，若 \(\displaystyle A\) 的伴随矩阵的秩等于1，则必有（　）。


（4）（续）（A）\(\displaystyle a=b\) 或 \(\displaystyle a+2b=0\)。（B）\(\displaystyle a=b\) 或 \(\displaystyle a+2b\ne0\)。（C）\(\displaystyle a\ne b\) 且 \(\displaystyle a+2b=0\)。（D）\(\displaystyle a\ne b\) 且 \(\displaystyle a+2b\ne0\)。


（5）设 \(\displaystyle \alpha_1,\allowbreak \alpha_2,\allowbreak \ldots,\allowbreak \alpha_s\) 均为 \(\displaystyle n\) 维向量，下列结论不正确的是（　）。
（A）若对于任意一组不全为零的数 \(\displaystyle k_1,\allowbreak \ldots,\allowbreak k_s\)，都有 \(\displaystyle k_1\alpha_1+\cdots+k_s\alpha_s\ne0\)，则向量组线性无关。（B）若向量组线性相关，则对于任意一组不全为零的数 \(\displaystyle k_1,\allowbreak \ldots,\allowbreak k_s\)，有 \(\displaystyle k_1\alpha_1+\cdots+k_s\alpha_s=0\)。（C）向量组线性无关的充分必要条件是此向量组的秩为 \(\displaystyle s\)。（D）向量组线性无关的必要条件是其中任意两个向量线性无关。


（6）将一枚硬币独立地掷两次，引进事件：\(\displaystyle A_1=\{\text{掷第一次出现正面}\}\)，\(\displaystyle A_2=\{\text{掷第二次出现正面}\}\)，\(\displaystyle A_3=\{\text{正、反面各出现一次}\}\)，\(\displaystyle A_4=\{\text{正面出现两次}\}\)，则事件（　）。
（A）\(\displaystyle A_1,\allowbreak A_2,\allowbreak A_3\) 相互独立。（B）\(\displaystyle A_2,\allowbreak A_3,\allowbreak A_4\) 相互独立。（C）\(\displaystyle A_1,\allowbreak A_2,\allowbreak A_3\) 两两独立。（D）\(\displaystyle A_2,\allowbreak A_3,\allowbreak A_4\) 两两独立。


\subsection*{三、（本题满分8分）}

设 \(\displaystyle f(x)=\dfrac{\displaystyle 1}{\displaystyle \pi x}+\dfrac{\displaystyle 1}{\displaystyle \sin\pi x}-\dfrac{\displaystyle 1}{\displaystyle \pi(1-x)}\)，\(\displaystyle x\in[\dfrac{\displaystyle 1}{\displaystyle 2},\allowbreak 1)\)，试补充定义 \(\displaystyle f(1)\) 使得 \(\displaystyle f(x)\) 在 \(\displaystyle [\dfrac{\displaystyle 1}{\displaystyle 2},\allowbreak 1]\) 上连续。


\subsection*{四、（本题满分8分）}

设 \(\displaystyle f(u,\allowbreak v)\) 具有二阶连续偏导数，且满足 \(\displaystyle \dfrac{\displaystyle \partial^2f}{\displaystyle \partial u^2}+\dfrac{\displaystyle \partial^2f}{\displaystyle \partial v^2}=1\)，又 \(\displaystyle g(x,\allowbreak y)=f[xy,\allowbreak \dfrac{\displaystyle 1}{\displaystyle 2}(x^2-y^2)]\)，求 \(\displaystyle \dfrac{\displaystyle \partial^2g}{\displaystyle \partial x^2}+\dfrac{\displaystyle \partial^2g}{\displaystyle \partial y^2}\)。


\subsection*{五、（本题满分8分）}

计算二重积分 \(\displaystyle I=\displaystyle \iint_D e^{-(x^2+y^2-\pi)}\sin(x^2+y^2)\,\mathrm dx\mathrm dy\)，其中积分区域 \(\displaystyle D=\{(x,\allowbreak y)\mid x^2+y^2\le\pi\}\)。


\subsection*{六、（本题满分9分）}

求幂级数 \(\displaystyle 1+\displaystyle \sum_{n=1}^{\infty}\dfrac{\displaystyle (-1)^nx^{2n}}{\displaystyle 2n}\;(|x|<1)\) 的和函数 \(\displaystyle f(x)\) 及其极值。


\subsection*{七、（本题满分9分）}

设 \(\displaystyle F(x)=f(x)g(x)\)，其中 \(\displaystyle f(x),\allowbreak g(x)\) 在 \(\displaystyle (-\infty,\allowbreak +\infty)\) 内满足以下条件：\(\displaystyle f^{\prime}(x)=g(x)\)，\(\displaystyle g^{\prime}(x)=f(x)\) 且 \(\displaystyle f(0)=0\)，\(\displaystyle f(x)+g(x)=2e^x\)。（1）求 \(\displaystyle F(x)\) 所满足的一阶微分方程；（2）求出 \(\displaystyle F(x)\) 的表达式。


\subsection*{八、（本题满分8分）}

设函数 \(\displaystyle f(x)\) 在 \(\displaystyle [0,\allowbreak 3]\) 上连续，在 \(\displaystyle (0,\allowbreak 3)\) 内可导，且 \(\displaystyle f(0)+f(1)+f(2)=3\)，\(\displaystyle f(3)=1\)。试证必存在 \(\displaystyle \xi\in(0,\allowbreak 3)\)，使 \(\displaystyle f^{\prime}(\xi)=0\)。


\subsection*{九、（本题满分13分）}

已知齐次线性方程组 \(\displaystyle \begin{cases}(a_1+b)x_1+a_2x_2+\cdots+a_nx_n=0,\allowbreak \\a_1x_1+(a_2+b)x_2+\cdots+a_nx_n=0,\allowbreak \\\cdots,\allowbreak \\a_1x_1+a_2x_2+\cdots+(a_n+b)x_n=0.\end{cases}\) 其中 \(\displaystyle \sum_{i=1}^na_i\ne0\)，试讨论 \(\displaystyle a_1,\allowbreak a_2,\allowbreak \ldots,\allowbreak a_n\) 和 \(\displaystyle b\) 满足何种关系时，（1）方程组仅有零解；（2）方程组有非零解，在有非零解时，求此方程组的一个基础解系。


\subsection*{十、（本题满分13分）}

设二次型 \(\displaystyle f(x_1,\allowbreak x_2,\allowbreak x_3)=X^{\mathrm T}AX=ax_1^2+2x_2^2-2x_3^2+2bx_1x_3\;(b>0)\)，其中二次型的矩阵 \(\displaystyle A\) 的特征值之和为1，特征值之积为 \(\displaystyle -12\)。（1）求 \(\displaystyle a,\allowbreak b\) 之值；（2）利用正交变换将二次型化为标准形，并写出所用的正交变换和对应的正交矩阵。


\subsection*{十一、（本题满分13分）}

设随机变量 \(\displaystyle X\) 的概率密度为 \(\displaystyle f(x)=\begin{cases}\dfrac{\displaystyle 1}{\displaystyle 3\sqrt[3]{x^2}},\allowbreak &x\in[1,\allowbreak 8],\allowbreak \\0,\allowbreak &\text{其他},\allowbreak \end{cases}\) \(\displaystyle F(x)\) 是 \(\displaystyle X\) 的分布函数。求随机变量 \(\displaystyle Y=F(X)\) 的分布函数。


\subsection*{十二、（本题满分13分）}

设随机变量 \(\displaystyle X\) 与 \(\displaystyle Y\) 独立，其中 \(\displaystyle X\) 的概率分布为 \(\displaystyle \begin{array}{c|cc}X&1&2\\\hline P&0.3&0.7\end{array}\)，而 \(\displaystyle Y\) 的概率密度为 \(\displaystyle f(y)\)，求随机变量 \(\displaystyle U=X+Y\) 的概率密度 \(\displaystyle g(u)\)。

\end{document}
